\input{TAC-801-SHARED-PREAMBLE.tex}
\begin{document}
\nolinenumbers
\begin{singlespace}
Franciscus Dylan Rosario, \\
7624 Melody \\
Rohnert Park, CA 94928 \\
E: soltrinox@gmail.com \\
Plaintiff, In Pro Per \\
\end{singlespace}
\vspace*{9mm}
\begin{tightcenter}
\textbf{SUPERIOR COURT OF THE STATE OF CALIFORNIA} \\
\textbf{COUNTY OF SAN FRANCISCO}
\end{tightcenter}
\vspace*{2.25mm}
\nolinenumbers
\begin{minipage}[t]{3in}
    \raggedright
    \textbf{Franciscus Dylan Rosario,} \\ \textbf{PLAINTIFF}, \\
    \hspace{1cm} v. \\
    \small
    CSAA Insurance Exchange; \\
    Carbone, Smith \& Koyama, LLP; \\
    Michael R. Chambers; \\
    Phillips, Spallas \& Angstadt LLP; \\
    Priya D. Navaratnasingham; \\
    Alberto Reyna; \\
    and DOES 1 through 20, inclusive; \\
    \normalsize
    \textbf{DEFENDANTS}
    \makebox[3in]{\hrulefill}
\end{minipage}%
\hspace{3mm}
\begin{minipage}[t]{0.5pt}
    \vspace{0pt}
    \rule{0.5pt}{3.1in}
\end{minipage}
\hspace{3mm}
\begin{minipage}[t]{3in}
    \raggedright
    \noindent \textbf{Case No.: CGC-25-631801} \\
    \textbf{(Related: CGC-21-594102)} \\
    ~\\
    \textbf{THIRD AMENDED COMPLAINT} \\
    \textbf{FOR EXTRINSIC FRAUD AND EQUITABLE RELIEF}
\end{minipage}
\vspace*{6mm}
\linenumbers
\begin{center}\textbf{TABLE OF CONTENTS}\end{center}
\vspace{2mm}
\CourtTOCRow{}{\tacctoclink{section-i-character-of-the-action}{SECTION I: CHARACTER OF THE ACTION}}{\pageref{toc:section-i-character-of-the-action}}
\CourtTOCRow{}{\tacctoclink{section-ii-noncommunicative-conduct-block}{SECTION II: NON-COMMUNICATIVE CONDUCT BLOCK}}{\pageref{toc:section-ii-noncommunicative-conduct-block}}
\CourtTOCRow{}{\tacctoclink{section-iii-standing-block}{SECTION III: STANDING BLOCK}}{\pageref{toc:section-iii-standing-block}}
\CourtTOCRow{}{\tacctoclink{section-iv-primary-rights-statement}{SECTION IV: PRIMARY RIGHTS STATEMENT}}{\pageref{toc:section-iv-primary-rights-statement}}
\CourtTOCRow{}{\tacctoclink{section-iv5-appeal-nonpreclusion}{SECTION IV.5: APPEAL NON-PRECLUSION}}{\pageref{toc:section-iv5-appeal-nonpreclusion}}
\CourtTOCRow{}{\tacctoclink{section-v-reservation-paragraph}{SECTION V: RESERVATION PARAGRAPH}}{\pageref{toc:section-v-reservation-paragraph}}
\CourtTOCRow{}{\tacctoclink{section-vi-tiered-prayer-for-relief}{SECTION VI: TIERED PRAYER FOR RELIEF}}{\pageref{toc:section-vi-tiered-prayer-for-relief}}
\CourtTOCRow{}{\tacctoclink{section-vii-partyspecific-allegations}{SECTION VII: PARTY-SPECIFIC ALLEGATIONS}}{\pageref{toc:section-vii-partyspecific-allegations}}
\CourtTOCRow{}{\tacctoclink{section-viii-day7-admission-foundation}{SECTION VIII: DAY-7 ADMISSION FOUNDATION}}{\pageref{toc:section-viii-day7-admission-foundation}}
\CourtTOCRow{}{\tacctoclink{section-ix-lawsbroken-catalog}{SECTION IX: LAWS-BROKEN CATALOG}}{\pageref{toc:section-ix-lawsbroken-catalog}}
\CourtTOCRow{}{\tacctoclink{causes-of-action}{CAUSES OF ACTION}}{\pageref{toc:causes-of-action}}
\CourtTOCRow{}{\tacctoclink{prayer-for-relief}{PRAYER FOR RELIEF}}{\pageref{toc:prayer-for-relief}}
\CourtTOCRow{}{\tacctoclink{verification}{VERIFICATION}}{\pageref{toc:verification}}
\vspace{4mm}
\phantomsection\label{toc:section-i-character-of-the-action}
\section{SECTION I: CHARACTER OF THE ACTION}\label{section-i-character-of-the-action}

\subsection{Independent Action in Equity Under Kougasian v. TMSL / Olivera v. Grace}\label{independent-action-in-equity-under-kougasian-v.-tmsl-olivera-v.-grace}

\begin{enumerate}
\def\labelenumi{\arabic{enumi}.}
\item
  This is an independent action in equity to set aside a judgment procured by extrinsic fraud, brought pursuant to the inherent equitable powers of this Court as recognized in \emph{Kougasian v. TMSL, Inc.} (2004) 118 Cal.App.4th 1028, 1032--35, and \emph{Olivera v. Grace} (1942) 19 Cal.2d 570, 575--76. This Court possesses inherent equity jurisdiction over actions to vacate judgments obtained through fraud on the tribunal. \emph{In re Marriage of Stevenot} (1984) 154 Cal.App.3d 1051.
\item
  This action is not an appeal, a motion for new trial, or a collateral attack on an otherwise valid judgment. It is a separate, original action invoking the Court's equity jurisdiction to vacate a judgment where the prevailing party's misconduct prevented full and fair adjudication of Plaintiff's claims.
\end{enumerate}

\subsection{Four-Channel Kulchar Framework}\label{four-channel-kulchar-framework}

\begin{enumerate}
\def\labelenumi{\arabic{enumi}.}
\setcounter{enumi}{2}
\item
  Plaintiff pleads extrinsic fraud under the four-channel framework of \emph{Kulchar v. Kulchar} (1969) 1 Cal.3d 467:

  \textbf{(a) Concealment of Facts:} Defendants fraudulently concealed the production history of the 911 audio, body-worn camera footage, and CAD logs, representing to the Court that these materials were ``not produced in discovery'' and of ``unknown provenance'' when they had been in defense counsel's possession since August-September 2022.

  \textbf{(b) False Promises / Representations of Non-Receipt:} Defense counsel made categorical representations to the tribunal in 2025 they ``never received'' critical evidence, when the certified record proves these materials were served on prior defense counsel in 2022 and transmitted during in 2023 file transfer.

  \textbf{(c) Keeping Plaintiff in Ignorance:} By maintaining 893 days of silence regarding their receipt of these materials (August 18, 2022 through January 27, 2025), Defendants kept Plaintiff in ignorance of the true state of the evidentiary record and induced Plaintiff to believe authentication was uncontested.

  \textbf{(d) Misleading the Tribunal:} Defendants procured structural evidentiary rulings by misleading the Court about their possession and the provenance of materials. The Court expressly articulated its reliance on defense counsel's officer-of-the-court status: ``I'm not going to question a lawyer's remarks regarding what they have and don't have. They are an officer of the Court.'' (Hr'g Tr. Feb.~28, 2025, pp.~291:26--292:7.)
\end{enumerate}

3.5. The channel of concealment pleaded here is the species of extrinsic fraud recognized by the California Supreme Court in \emph{In re Marriage of Modnick} (1983) 33 Cal.3d 897, 907, where the Court held that a party who conceals material matter from the tribunal does not merely commit a litigation misstep---he effects a ``denial of the very existence'' of the concealed matter, with the doctrinal consequence that the subject is ``never fairly before the court for adjudication.'' (\emph{Id.} at p.~907, quoting \emph{Clark v. Clark} (1961) 195 Cal.App.2d 373, 381; accord \emph{Jorgensen v. Jorgensen} (1948) 32 Cal.2d 13, 21.) \emph{Modnick} thus locates the wrong not in the falsity of any particular representation (which would be intrinsic) but in the structural \emph{suppression of subject matter} from the adversary process itself. That is precisely the wrong pleaded here.

\textbf{3.5.1. The Suppressed Subject Matter.} The materials Defendants concealed for 893 days---the 911 audio of August 18, 2022, the San Francisco Police Department body-worn camera footage, the Computer-Aided Dispatch (CAD) entries, and the supplemental police report---are not peripheral impeachment. They are the primary, contemporaneous, government-generated record of the very incident on which Defendants' defense, their declarations, and the Court's structural evidentiary rulings all depended. They \emph{constitute} the evidentiary record of August 18, 2022.

\textbf{3.5.2. The Denial.} Defendants did not merely dispute the weight of these materials; through counsel, and through sustained silence coupled with affirmative officer-of-the-court representations, they denied that the materials had ever been received. (¶¶ 3.3--3.4 supra; Hr'g Tr. Feb.~28, 2025, pp.~291:26--292:7; MIL No.~17 (Jan.~17, 2025) (``no idea where this 911 call log or audio recording came from'').) Under \emph{Modnick}, that denial is not a credibility dispute to be resolved at trial. It is a denial of the existence of the very evidentiary record the Court was being asked to adjudicate---precisely the structural operation \emph{Modnick} identifies as extrinsic fraud.

\textbf{3.5.3. The Doctrinal Consequence.} Because the concealed subject matter was ``never fairly before the court for adjudication'' (\emph{Modnick}, supra, 33 Cal.3d at p.~907), every ruling predicated on the assumed completeness of the August 18, 2022 record---including the authentication ruling, the seven consequential rulings catalogued at ¶ 5 infra, and every case-management ruling crediting defense counsel's officer-of-the-court representations---was procured in a proceeding from which the decisive evidentiary subject matter had been structurally withheld. That is extrinsic fraud \emph{as a matter of California Supreme Court doctrine}, not appellate gloss, and it is the same doctrinal move the Supreme Court has since reaffirmed in \emph{Kulchar v. Kulchar} (1969) 1 Cal.3d 467, 471, and \emph{Estate of Sanders} (1985) 40 Cal.3d 607, 614--616.

\textbf{3.5.4. Why the Intrinsic-Fraud Line Does Not Reach This Pleading.} The rule of \emph{Pico v. Cohn} (1891) 91 Cal. 129, 133--134---that perjury alone is intrinsic---is not implicated here and cannot be used to re-characterize this pleading. Plaintiff does not allege that Defendants testified falsely about the \emph{contents} of the 911 audio, body-worn camera footage, CAD entries, or supplemental report. Plaintiff alleges that Defendants denied the \emph{existence} of those materials in their possession, and by that denial prevented the tribunal from ever confronting them. \emph{Modnick} itself drew exactly that line (33 Cal.3d at pp.~905--907), holding that concealment-as-denial-of-existence is categorically distinct from perjury-about-a-matter-in-issue and falls on the extrinsic side of \emph{Pico}. The pleading here tracks that distinction on its face.

\subsection{\textsection{}I.4.5 --- Supreme Court--Recognized Right to a Fair Adversary Hearing}\label{i.4.5-supreme-courtrecognized-right-to-a-fair-adversary-hearing}

3.6. The primary right invoked in this action is the right recognized by the California Supreme Court in \emph{Westphal v. Westphal} (1942) 20 Cal.2d 393, 397: freedom from a judgment procured by ``fraud {[}that{]} is extrinsic \ldots{} where the unsuccessful party has been prevented from exhibiting fully his case, by fraud or deception practiced on him by his opponent.'' That formulation has been reaffirmed by the Supreme Court in \emph{Kulchar v. Kulchar} (1969) 1 Cal.3d 467, 471 (``a fair adversary hearing'' is the ``primary right'' protected), \emph{In re Marriage of Modnick} (1983) 33 Cal.3d 897, 905 (extrinsic fraud is a ``broad concept'' reaching ``any extrinsic circumstance'' that has ``prevented a party from fully participating''), and \emph{Estate of Sanders} (1985) 40 Cal.3d 607, 614--616 (fiduciary breach by an officer of the court ``effectively prevented'' the opposing party from ``a fair adversary hearing,'' quoting \emph{Larrabee v. Tracy} (1943) 21 Cal.2d 645, 650--651).

3.7. The injury pleaded in this action is therefore measured by what Plaintiff was \emph{prevented} from presenting, not by what eventually leaked into the record under a corrupted procedural posture. Under \emph{Westphal}, \emph{Kulchar}, \emph{Modnick}, and \emph{Sanders}, the Supreme Court has located the operative harm at the point of deprivation of opportunity, not at the point of evidentiary admission. That is the primary right this independent action in equity vindicates.

\subsection{\textsection{}I.5 --- The Court's Articulated Reliance and Consequential Rulings}\label{i.5-the-courts-articulated-reliance-and-consequential-rulings}

\begin{enumerate}
\def\labelenumi{\arabic{enumi}.}
\setcounter{enumi}{3}
\tightlist
\item
  On February 28, 2025, the Court articulated the institutional basis for accepting defense counsel's representations:
\end{enumerate}

\begin{quote}
``I'm not going to question a lawyer's remarks regarding what they have and don't have. They are an officer of the Court, and they are supposed to conduct themselves appropriately before the Court. If they are not, there are other sanctions that they could suffer here. So far I haven't heard or seen anything on the part of the current law firm which tells me that they have conducted themselves improperly.'' --- \emph{The Court (Garrett L. Wong), Hr'g Tr. Feb.~28, 2025, pp.~291:26--292:7.}
\end{quote}

\begin{enumerate}
\def\labelenumi{\arabic{enumi}.}
\setcounter{enumi}{4}
\item
  Based on this misplaced reliance on defense counsel's representations of non-receipt, the Court issued the following consequential rulings:

  \begin{enumerate}
  \def\labelenumii{(\alph{enumii})}
  \item
    \textbf{MIL No.~1 Denied} --- Jan.~27, 2025, p.~41:3--11 (excluding 911 materials);
  \item
    \textbf{911 CDs Excluded} --- Jan.~27, 2025, p.~38:22--25;
  \item
    \textbf{911 Transcripts ``Never Admitted''} --- Jan.~27, 2025, p.~39:18--24 (categorical exclusion);
  \item
    \textbf{Continuance Denied} --- Feb.~18, 2025, p.~181;
  \item
    \textbf{\textsection{} 473 ``Surprise'' Relief Rejected} --- Feb.~18, 2025, p.~181:20--25;
  \item
    \textbf{Burden-Shift Imposed} --- Feb.~28, 2025, p.~293:1--3 (``it's required of you to send the documents over. Otherwise, I won't allow them in.'');
  \item
    \textbf{Supplemental Report Excluded} --- RT 7, p.~156:24--28.
  \end{enumerate}
\item
  Each of these rulings deprived Plaintiff of the ability to present critical evidence at trial. The deprivation was caused by defense counsel's misrepresentations about their custody, receipt, and knowledge of the materials.
\end{enumerate}

\subsection{Gravamen: Non-Communicative Procurement}\label{gravamen-non-communicative-procurement}

\begin{enumerate}
\def\labelenumi{\arabic{enumi}.}
\setcounter{enumi}{6}
\item
  The gravamen of this action is extrinsic fraud within the meaning of the California Supreme Court's governing sentence in \emph{Westphal v. Westphal} (1942) 20 Cal.2d 393, 397 --- fraud that ``prevented {[}the unsuccessful party{]} from exhibiting fully his case'' --- and within the ``broad concept'' of extrinsic fraud restated in \emph{Kulchar v. Kulchar} (1969) 1 Cal.3d 467, 471, \emph{In re Marriage of Modnick} (1983) 33 Cal.3d 897, 905, and \emph{Estate of Sanders} (1985) 40 Cal.3d 607, 614--616. Defendants procured the underlying judgment through a combination of non-communicative acts and institutional breaches that, under \emph{Sanders}, 40 Cal.3d at 615, convert otherwise in-court conduct into extrinsic fraud because it was carried out by officers of the court owing a fiduciary duty of candor to the tribunal:

  \begin{enumerate}
  \def\labelenumii{(\alph{enumii})}
  \item
    Receiving materials in discovery in August-September 2022 and proceeding through sworn deposition course without timely written authenticity challenge;
  \item
    Maintaining 893 days of silence regarding receipt;
  \item
    Executing the May 4, 2023 file transfer without conveying the production history to successor counsel;
  \item
    Filing court-facing representations that reversed the prior out-of-court posture; and
  \item
    Directing and supervising the defense file in a manner that ensured concealment.
  \end{enumerate}
\end{enumerate}

7.5. Under \emph{In re Marriage of Modnick} (1983) 33 Cal.3d 897, 907, this conduct produced a record in which the concealed materials and their production history were ``never fairly before the court for adjudication.'' The Court's seven consequential rulings (¶ 5 supra) were issued in a posture in which the existence, provenance, and chain-of-custody of the 911 audio, BWC footage, CAD logs, and supplemental report had been denied by officers of the court. That is the paradigmatic \emph{Modnick} defect. It is extrinsic as a matter of binding California Supreme Court doctrine.

\begin{center}\rule{0.5\linewidth}{0.5pt}\end{center}

\phantomsection\label{toc:section-ii-noncommunicative-conduct-block}
\section{SECTION II: NON-COMMUNICATIVE CONDUCT BLOCK}\label{section-ii-non-communicative-conduct-block}

\subsection{Privilege Carveout --- Civil Code \textsection{} 47(b) Inapplicable}\label{privilege-carveout-civil-code-47b-inapplicable}

\begin{enumerate}
\def\labelenumi{\arabic{enumi}.}
\setcounter{enumi}{7}
\item
  The conduct complained of herein consists entirely of non-communicative acts that fall outside the protection of Civil Code \textsection{} 47(b):

  \textbf{(a) Years-Long Silence:} Defense counsel's 893 days of silence from August 18, 2022 through January 27, 2025 regarding receipt of the 911 audio, body-worn camera footage, and CAD logs is non-communicative conduct.

  \textbf{(b) File-Transfer Concealment:} The May 4, 2023 institutional file transfer from Chambers/CSK to successor counsel PSA without conveying the production history is non-communicative conduct---an act of non-disclosure, not a statement.

  \textbf{(c) Custodial Possession:} Maintaining custody of discovery materials while representing ignorance of their provenance is conduct, not communication.

  \textbf{(d) Failure to Correct:} The failure to correct materially false representations after the April 16, 2025 admission that materials were ``provided by Dolan'' is non-communicative conduct.
\end{enumerate}

\subsection{The Nine-Point Evidentiary Chain of Defense Receipt}\label{the-nine-point-evidentiary-chain-of-defense-receipt}

\begin{enumerate}
\def\labelenumi{\arabic{enumi}.}
\setcounter{enumi}{8}
\tightlist
\item
  The denials of receipt were not merely mistaken; they were irreconcilable with a documented two-year record of notice and possession:
\end{enumerate}

{\def\LTcaptype{table} % do not increment counter
\begin{longtable}[]{@{}
  >{\raggedright\arraybackslash}p{(\linewidth - 8\tabcolsep) * \real{0.2000}}
  >{\raggedright\arraybackslash}p{(\linewidth - 8\tabcolsep) * \real{0.2000}}
  >{\raggedright\arraybackslash}p{(\linewidth - 8\tabcolsep) * \real{0.2000}}
  >{\raggedright\arraybackslash}p{(\linewidth - 8\tabcolsep) * \real{0.2000}}
  >{\raggedright\arraybackslash}p{(\linewidth - 8\tabcolsep) * \real{0.2000}}@{}}
\toprule\noalign{}
\begin{minipage}[b]{\linewidth}\raggedright
\#
\end{minipage} & \begin{minipage}[b]{\linewidth}\raggedright
Date
\end{minipage} & \begin{minipage}[b]{\linewidth}\raggedright
Proof Artifact
\end{minipage} & \begin{minipage}[b]{\linewidth}\raggedright
Pin Cite / Source
\end{minipage} & \begin{minipage}[b]{\linewidth}\raggedright
Who Had Notice
\end{minipage} \\
\midrule\noalign{}
\endhead
\bottomrule\noalign{}
\endlastfoot
1 & Aug.~18, 2022 & Service of 911 audio by Dolan to defense counsel & Jessup testimony, Hr'g Tr. Feb.~18, 2025, pp.~139--140 & Chambers / CSK \\
2 & Aug.~18, 2022 & Rosario deposition --- Chambers present & Rosario depo tr. & Chambers \\
3 & Sept.~30, 2022 & Abdelhalim deposition --- Chambers present; 911-related sworn testimony & Abdelhalim depo tr. & Chambers \\
4 & May 4, 2023 & Substitution of Attorney (Transaction ID \#69953800) --- file transferred from CSK {\textrightarrow{}} PSA & CT record & CSK {\textrightarrow{}} PSA \\
5 & Feb.~1, 2024 & Reyna (PSA) Compex subpoena to SFPD for identical materials & Order No.~CA1147102-021 & PSA \\
6 & Feb.~26, 2024 & SFPD produced records responsive to subpoena & SFPD Report 210079556 & PSA \\
7 & Sept.~9, 2024 & Plaintiff's FROG/SROG served on Reyna identifying CAD log & FROG p.~28, ll. 23--25; SROG p.~9, ll. 9--11 & PSA \\
8 & Feb.~18, 2025 & Jessup testifies ``Dolan had a copy since 2021'' & Hr'g Tr. Feb.~18, 2025, p.~139, l. 11 & Navaratnasingham (present) \\
9 & Apr.~16, 2025 & Navaratnasingham admission: ``They were provided by Dolan'' & RT 7, p.~213, ll. 13--15 & PSA (admitting) \\
\end{longtable}
}

\subsection{Flatley Illegality Anchor}\label{flatley-illegality-anchor}

\begin{enumerate}
\def\labelenumi{\arabic{enumi}.}
\setcounter{enumi}{9}
\item
  Defendants' conduct constitutes illegal conduct as a matter of law under Business and Professions Code \textsection{} 6128(a), which provides: ``Every attorney is guilty of a misdemeanor who\ldots{} {[}i{]}s guilty of any deceit or collusion, or consents to any deceit or collusion, with intent to deceive the court or any party.''
\item
  Because Defendants' conduct falls within \textsection{} 6128(a), it is not protected by the litigation privilege under \emph{Flatley v. Mauro} (2006) 39 Cal.4th 299, 320--325 (conduct that violates a statute constituting ``criminal'' conduct is not protected activity).
\item
  Additionally, \emph{Action Apartment Ass'n, Inc.~v. City of Santa Monica} (2007) 41 Cal.4th 1232, 1248--49, establishes that extrinsic fraud in the procurement of a judgment falls outside the scope of the litigation privilege. The privilege is designed to facilitate open advocacy in judicial proceedings, not to shield conduct that subverts the judicial process itself.
\end{enumerate}

\subsection{Pre-Rebut Throckmorton}\label{pre-rebut-throckmorton}

\begin{enumerate}
\def\labelenumi{\arabic{enumi}.}
\setcounter{enumi}{12}
\tightlist
\item
  To the extent Defendants invoke the federal ``intrinsic/extrinsic'' ceiling of \emph{United States v. Throckmorton} (1878) 98 U.S. 61, that argument is foreclosed by three controlling California Supreme Court decisions. \emph{Kulchar v. Kulchar} (1969) 1 Cal.3d 467, 471, holds that extrinsic fraud exists whenever a party has been deprived of ``a fair adversary hearing.'' \emph{In re Marriage of Modnick} (1983) 33 Cal.3d 897, 905, describes extrinsic fraud as a ``broad concept'' reaching ``any extrinsic circumstance'' that has ``prevented a party from fully participating'' in the proceeding. \emph{Estate of Sanders} (1985) 40 Cal.3d 607, 614--616, holds that a fiduciary breach by an officer of the court that ``effectively prevented'' a fair adversary hearing (quoting \emph{Larrabee v. Tracy} (1943) 21 Cal.2d 645, 650--651) is extrinsic fraud as a matter of state Supreme Court doctrine, notwithstanding that the breach occurred through representations made in court. \emph{Throckmorton} is not binding on California courts, is not California Supreme Court authority, and cannot override \emph{Westphal}, \emph{Kulchar}, \emph{Modnick}, or \emph{Sanders}.
\end{enumerate}

\begin{center}\rule{0.5\linewidth}{0.5pt}\end{center}

\phantomsection\label{toc:section-iii-standing-block}
\section{SECTION III: STANDING BLOCK}\label{section-iii-standing-block}

\subsection{Kougasian/Olivera Equity-Action Standing}\label{kougasianolivera-equity-action-standing}

\begin{enumerate}
\def\labelenumi{\arabic{enumi}.}
\setcounter{enumi}{13}
\tightlist
\item
  Plaintiff has standing to bring this independent action in equity because:
\end{enumerate}

\begin{enumerate}
\def\labelenumi{(\alph{enumi})}
\item
  Plaintiff was a party to the underlying litigation (CGC-21-594102);
\item
  A judgment was entered against Plaintiff in that litigation on April 25, 2025;
\item
  Plaintiff alleges the judgment was procured by extrinsic fraud; and
\item
  Plaintiff has no adequate remedy at law because the ordinary appellate process cannot address evidence that was fraudulently concealed and institutional misconduct that operated outside the trial itself.
\end{enumerate}

\subsection{Shaolian/Norton/Estate of Sanders Distinction}\label{shaoliannortonestate-of-sanders-distinction}

\begin{enumerate}
\def\labelenumi{\arabic{enumi}.}
\setcounter{enumi}{14}
\tightlist
\item
  Defendants may cite \emph{Shaolian v. Safeco Ins. Co.} (1999) 71 Cal.App.4th 268 and \emph{Norton v. Hines} (1975) 49 Cal.App.3d 917 for the proposition that a non-client cannot sue an attorney for litigation conduct. However:
\end{enumerate}

\begin{enumerate}
\def\labelenumi{(\alph{enumi})}
\item
  \emph{Shaolian} addresses direct actions for policy benefits against insurers; it does not address equitable fraud-on-the-court theories brought by a party to the underlying litigation seeking vacatur of a judgment procured by fraud.
\item
  \emph{Norton} concerns duty owed to non-clients in ordinary malpractice scenarios, not independent actions in equity to vacate judgments.
\item
  \emph{Estate of Sanders} (1985) 40 Cal.3d 607, 615--616 establishes that attorneys who participate in fraud on the court are proper parties to an action for equitable relief. Counsel who misrepresent material facts to the tribunal waive any privilege that would otherwise shield them from suit by the party against whom the fraud was perpetrated.
\end{enumerate}

15.5. The Attorney Defendants occupied the status of officers of the court, a status the California Supreme Court in \emph{Estate of Sanders} (1985) 40 Cal.3d 607, 615, treats as imposing a fiduciary duty that runs to the integrity of the proceeding and, derivatively, to the opposing party's primary right to a fair adversary hearing. The February 28, 2025 colloquy (Hr'g Tr. pp.~291:26--292:7) is the on-record predicate that the trial court extended institutional reliance to that fiduciary status. Under \emph{Sanders}, 40 Cal.3d at 615--616, an officer-of-the-court fiduciary breach that ``effectively prevented'' the opposing party from a fair adversary hearing is extrinsic fraud as a matter of California Supreme Court doctrine, and the opposing party has standing in equity to sue for its cure --- regardless of the absence of an attorney-client relationship between the party and counsel. This paragraph states that standing theory in express \emph{Sanders} terms.

\subsection{Agency/Ratification Under Civil Code \textsection{} 2338}\label{agencyratification-under-civil-code-2338}

\begin{enumerate}
\def\labelenumi{\arabic{enumi}.}
\setcounter{enumi}{15}
\item
  CSAA Insurance Exchange is liable for the acts of its agents, Defendants CSK, PSA, Chambers, Navaratnasingham, and Reyna, under Civil Code \textsection{} 2338, which provides that a principal is responsible for the negligent and wrongful acts of its agent in the transaction of the business of the agency.
\item
  CSAA Insurance Exchange ratified the conduct of its agents by:
\end{enumerate}

\begin{enumerate}
\def\labelenumi{(\alph{enumi})}
\item
  Paying attorney's fees for services that included the concealment of production history;
\item
  Directing litigation strategy through the insurer-defense relationship;
\item
  Approving the May 4, 2023 substitution of counsel;
\item
  Accepting the benefit of the judgment procured through concealment; and
\item
  Failing to repudiate counsel's conduct after learning of the April 16, 2025 admission.
\end{enumerate}

\begin{enumerate}
\def\labelenumi{\arabic{enumi}.}
\setcounter{enumi}{17}
\tightlist
\item
  Under the \emph{Doctors' Company v. Superior Court} (1989) 49 Cal.3d 39 aiding-and-abetting framework, an insurer that acts outside its normal insurer role by directing, approving, or ratifying fraudulent litigation conduct may be held directly liable for that conduct.
\end{enumerate}

\subsection{Business and Professions Code \textsection{} 6128(a) as Independent Tort}\label{business-and-professions-code-6128a-as-independent-tort}

\begin{enumerate}
\def\labelenumi{\arabic{enumi}.}
\setcounter{enumi}{18}
\tightlist
\item
  Under the \emph{Rickley v. Goodfriend} (2013) 212 Cal.App.4th 1136 and \emph{Shaffer v. Superior Court} (1995) 33 Cal.App.4th 993 line of cases, a violation of Business and Professions Code \textsection{} 6128(a) gives rise to an independent tort claim by the injured party. The statute creates a private right of action with treble damages.
\end{enumerate}

\subsection{Constructive Fraud Under Civil Code \textsection{} 1573}\label{constructive-fraud-under-civil-code-1573}

\begin{enumerate}
\def\labelenumi{\arabic{enumi}.}
\setcounter{enumi}{19}
\tightlist
\item
  Defendants' concealment of the production history and failure to correct false representations to the tribunal constitutes constructive fraud under Civil Code \textsection{} 1573, which defines constructive fraud as:
\end{enumerate}

\begin{quote}
``(1) In any breach of duty which, without an actually fraudulent intent, gains an advantage to the person in fault, or anyone claiming under him, by misleading another to his prejudice, or to the prejudice of anyone claiming under him; or

\begin{enumerate}
\def\labelenumi{(\arabic{enumi})}
\setcounter{enumi}{1}
\tightlist
\item
  In any such act or omission as the law specially declares to be fraudulent, without respect to actual fraud.''
\end{enumerate}
\end{quote}

\begin{enumerate}
\def\labelenumi{\arabic{enumi}.}
\setcounter{enumi}{20}
\tightlist
\item
  The February 28, 2025 colloquy establishes that the tribunal extended confidential institutional reliance to counsel's representations about possession and provenance. Civil Code \textsection{} 1573 supplies the civil-side vocabulary for breach of that kind of duty where the record supports misleading acts or nondisclosures. \emph{Assilzadeh v. California Federal Bank} (2000) 82 Cal.App.4th 399, 415.
\end{enumerate}

\begin{center}\rule{0.5\linewidth}{0.5pt}\end{center}

\phantomsection\label{toc:section-iv-primary-rights-statement}
\section{SECTION IV: PRIMARY RIGHTS STATEMENT}\label{section-iv-primary-rights-statement}

\subsection{Distinct Primary Right: Freedom from Fraud on the Court}\label{distinct-primary-right-freedom-from-fraud-on-the-court}

\begin{enumerate}
\def\labelenumi{\arabic{enumi}.}
\setcounter{enumi}{21}
\item
  The primary right asserted in this action is Plaintiff's right to be free from having a judgment procured against him through fraud on the court. Under \emph{Mycogen Corp.~v. Monsanto Co.} (2002) 28 Cal.4th 888, 904, and \emph{Boeken v. Philip Morris USA, Inc.} (2010) 48 Cal.4th 788, 798, the controlling formulation is that ``the primary right is simply the plaintiff's right to be free from the particular injury suffered.''
\item
  The primary right vindicated here is distinct from:
\end{enumerate}

\begin{enumerate}
\def\labelenumi{(\alph{enumi})}
\item
  The personal injury primary right litigated in CGC-21-594102 (right to be free from negligent bodily harm);
\item
  Any right to accurate claim handling asserted in CGC-25-631802 (right to be free from fraudulent claim denial); and
\item
  Any right to policy benefits (not asserted in either action).
\end{enumerate}

\subsection{Issues Not Adjudicated --- Five-Issue Block}\label{issues-not-adjudicated-five-issue-block}

\begin{enumerate}
\def\labelenumi{\arabic{enumi}.}
\setcounter{enumi}{23}
\tightlist
\item
  The following issues were \textbf{never presented to, litigated before, or necessarily decided by} the trier of fact in the underlying action:
\end{enumerate}

\begin{enumerate}
\def\labelenumi{(\alph{enumi})}
\item
  Whether the production history of the 911 audio, body-worn camera footage, and CAD records was concealed from successor counsel during the May 4, 2023 file transfer;
\item
  Whether prior defense counsel's out-of-court positions regarding authenticity of these materials were consistent with their internal records;
\item
  Whether CSAA directed or ratified the file-transfer concealment and the subsequent court-facing misrepresentations;
\item
  Whether the out-of-court positions taken by prior defense counsel contradicted the later tribunal-facing positions taken by successor counsel; and
\item
  Whether Plaintiff was deprived of a fair adversary hearing by extrinsic fraud operating outside the trial itself.
\end{enumerate}

\subsection{Primary Rights Comparison Table}\label{primary-rights-comparison-table}

\begin{enumerate}
\def\labelenumi{\arabic{enumi}.}
\setcounter{enumi}{24}
\tightlist
\item
  Under \emph{Mycogen} and \emph{Boeken}, the following comparison demonstrates that no ``same primary right'' match exists between the underlying action and this action:
\end{enumerate}

{\def\LTcaptype{table} % do not increment counter
\begin{longtable}[]{@{}
  >{\raggedright\arraybackslash}p{(\linewidth - 4\tabcolsep) * \real{0.1268}}
  >{\raggedright\arraybackslash}p{(\linewidth - 4\tabcolsep) * \real{0.4789}}
  >{\raggedright\arraybackslash}p{(\linewidth - 4\tabcolsep) * \real{0.3944}}@{}}
\toprule\noalign{}
\begin{minipage}[b]{\linewidth}\raggedright
Element
\end{minipage} & \begin{minipage}[b]{\linewidth}\raggedright
Underlying Trial (CGC-21-594102)
\end{minipage} & \begin{minipage}[b]{\linewidth}\raggedright
This Action (CGC-25-631801)
\end{minipage} \\
\midrule\noalign{}
\endhead
\bottomrule\noalign{}
\endlastfoot
\textbf{Primary Right} & Freedom from negligent bodily harm from the collision & Freedom from a judgment procured by extrinsic fraud; right to a fair adversary hearing \\
\textbf{Duty} & Ordinary care in operating a motor vehicle & Candor to the tribunal; conveyance of production history on substitution; institutional duties of principal and counsel \\
\textbf{Injury} & Physical and economic harm from the collision & Deprivation of a fair adversary hearing and attendant procedural harm \\
\end{longtable}
}

\begin{enumerate}
\def\labelenumi{\arabic{enumi}.}
\setcounter{enumi}{25}
\tightlist
\item
  Because there is no ``same primary right'' match, res judicata and collateral estoppel do not bar this complaint. \emph{Lucido v. Superior Court} (1990) 51 Cal.3d 335.
\end{enumerate}

26.5. The primary-rights analysis here is governed by four California Supreme Court decisions, each of which independently defeats any defense preclusion theory. \emph{Mycogen Corp.~v. Monsanto Co.} (2002) 28 Cal.4th 888, 904, holds that the primary right ``is simply the plaintiff's right to be free from the particular injury suffered.'' \emph{Boeken v. Philip Morris USA, Inc.} (2010) 48 Cal.4th 788, 798, holds that the primary right is independent of the legal theory or forum. \emph{Crowley v. Katleman} (1994) 8 Cal.4th 666, 681--682, holds that separate injuries give rise to separate primary rights. \emph{Lucido v. Superior Court} (1990) 51 Cal.3d 335, 341, holds that issue preclusion attaches only to issues ``actually litigated'' and ``necessarily decided.'' Under \emph{Mycogen/Boeken}, the injury here --- deprivation of a fair adversary hearing by officer-of-the-court fraud --- is a distinct injury from the bodily-harm injury litigated in CGC-21-594102, so it states a separate primary right. Under \emph{Crowley}, that separateness is not defeated by shared underlying facts. Under \emph{Lucido}, no issue decided in the underlying action was ``actually litigated'' on the question of defense counsel's concealment, because the very existence of the concealed record was hidden from the trier of fact. This TAC is not barred by any of the four Supreme Court preclusion doctrines the defense is expected to invoke.

\subsection{Mycogen/Crowley/Lucido Recital}\label{mycogencrowleylucido-recital}

\begin{enumerate}
\def\labelenumi{\arabic{enumi}.}
\setcounter{enumi}{26}
\tightlist
\item
  This action satisfies the requirements for an independent action in equity as set forth in \emph{Mycogen Corp.~v. Monsanto Co.} (2002) 28 Cal.4th 888, \emph{Crowley v. Katleman} (1994) 8 Cal.4th 666, and \emph{Lucido v. Superior Court} (1990) 51 Cal.3d 335:
\end{enumerate}

\begin{enumerate}
\def\labelenumi{(\alph{enumi})}
\item
  The fraud is extrinsic (relates to procurement, not merits);
\item
  Plaintiff had no opportunity to present a full and fair case due to concealment;
\item
  Plaintiff is without fault or negligence contributing to the fraud;
\item
  Plaintiff exercised diligence in discovering the fraud and seeking relief; and
\item
  No adequate remedy at law exists (appeal cannot address institutional misconduct and newly discovered evidence of fraud).
\end{enumerate}

27.5. Plaintiff's diligence is to be measured against the California Supreme Court's diligence standard in \emph{Kulchar v. Kulchar} (1969) 1 Cal.3d 467, 472--473, which holds that a party is not required to discover concealment uniquely within the other party's possession until a triggering disclosure occurs. Here, the triggering disclosure is the April 16, 2025 on-record admission by Defendant Navaratnasingham that the materials ``were provided by Dolan'' (RT 7, p.~213, ll. 13--15). That admission --- which is the first time defense counsel acknowledged the true provenance on the record --- is the \emph{Kulchar} discovery trigger. Plaintiff's investigation, reconstruction of the nine-point evidentiary chain, and filing of this action within the applicable limitations period is diligent as a matter of California Supreme Court doctrine.

\begin{center}\rule{0.5\linewidth}{0.5pt}\end{center}

\phantomsection\label{toc:section-iv5-appeal-nonpreclusion}
\section{SECTION IV.5: APPEAL NON-PRECLUSION}\label{section-iv.5-appeal-non-preclusion}

\begin{enumerate}
\def\labelenumi{\arabic{enumi}.}
\setcounter{enumi}{27}
\item
  Plaintiff has filed a notice of appeal in Case No.~A173827, which concerns instructional error in CGC-21-594102. That appeal addresses issues related to jury instructions and trial court rulings on evidentiary matters adjudicated during the trial itself.
\item
  This independent action concerns extrinsic fraud in the procurement of the judgment---specifically, institutional misconduct and concealment that operated \textbf{outside} the trial itself and could not have been raised in the underlying trial or on direct appeal.
\item
  Under \emph{Kougasian v. TMSL, Inc.} (2004) 118 Cal.App.4th 1028, 1034--36, and \emph{In re Marriage of Stevenot} (1984) 154 Cal.App.3d 1051, an independent equitable action is cognizable alongside a direct appeal because the two remedies address different wrongs. The appeal addresses rulings made during the trial; this action addresses the fraud that corrupted the institutional process leading to those rulings.
\item
  Defense reliance on federal \emph{United States v. Throckmorton} (1878) 98 U.S. 61 is unavailing because \emph{Throckmorton} is not binding California authority, and California has adopted the broader \emph{Kulchar/Kougasian/Olivera} framework recognizing independent equity actions for fraud that denies a party the opportunity to fully and fairly present his case.
\end{enumerate}

\begin{center}\rule{0.5\linewidth}{0.5pt}\end{center}

\phantomsection\label{toc:section-v-reservation-paragraph}
\section{SECTION V: RESERVATION PARAGRAPH}\label{section-v-reservation-paragraph}

\subsection{Express Disclaimers}\label{express-disclaimers}

\begin{enumerate}
\def\labelenumi{\arabic{enumi}.}
\setcounter{enumi}{31}
\tightlist
\item
  \textbf{Plaintiff expressly disclaims and does not seek in this action:}
\end{enumerate}

\begin{enumerate}
\def\labelenumi{(\alph{enumi})}
\item
  \textbf{Personal Injury Damages from CGC-21-594102:} Plaintiff does not seek to recover any personal injury damages that were or could have been awarded in CGC-21-594102. The personal injury primary right is not at issue in this action.
\item
  \textbf{Prevailing-Party Costs from Underlying Trial:} Plaintiff does not seek to recover costs awarded to Defendants as prevailing parties in CGC-21-594102, except to the extent such costs would be vacated as a consequence of vacating the underlying judgment.
\item
  \textbf{Duplicative 802 Damages:} Plaintiff does not seek to recover in this action any damages that are also sought in CGC-25-631802. To the extent any damages overlap, Plaintiff will elect remedies to avoid duplication.
\item
  \textbf{Pro Per Attorney's Fees:} Plaintiff, proceeding pro se, does not seek attorney's fees for self-representation in this action.
\end{enumerate}

\subsection{\textsection{}V.5 --- Partial-Presentation Rebuttal}\label{v.5-partial-presentation-rebuttal}

\begin{enumerate}
\def\labelenumi{\arabic{enumi}.}
\setcounter{enumi}{32}
\item
  Extrinsic fraud is measured by the deprivation of the \emph{opportunity to present} the case, not by what ultimately leaked into the record. \emph{Westphal v. Westphal} (1942) 20 Cal.2d 393, 397 (``prevented from exhibiting fully his case''); \emph{Kulchar v. Kulchar} (1969) 1 Cal.3d 467, 471 (``primary right'' is ``a fair adversary hearing''). Partial presentation under a corrupted procedural posture is not a cure under the Supreme Court's deprivation rule. The concealment of the production history altered what Plaintiff could authenticate, subpoena, and argue before and during trial --- i.e., it altered the opportunity itself, which is the metric the Supreme Court uses.
\item
  The fact that the 911 audio was ultimately admitted on Day 8 of trial---after Defendant Abdelhalim authenticated it as his own voice---does not cure the extrinsic fraud, because:
\end{enumerate}

\begin{enumerate}
\def\labelenumi{(\alph{enumi})}
\item
  Plaintiff was forced to testify before the 911 audio was admitted, disrupting his intended evidentiary sequence;
\item
  The categorical exclusion of the 911 transcripts (Jan.~27, 2025, p.~39:18--24) was never reversed;
\item
  The supplemental police report remained excluded;
\item
  Plaintiff was denied a reporter read-back of Abdelhalim's impeachment confirmation when defense argued in closing that Plaintiff ``changed his story'' (RT 10, pp.~49--52); and
\item
  The burden-shift ruling forced Plaintiff to re-produce materials already in defense's possession.
\end{enumerate}

34.5. The applicable standard is \emph{In re Marriage of Modnick} (1983) 33 Cal.3d 897, 907: the concealed materials and their production history must have been ``fairly before the court for adjudication.'' Under that standard, the seven consequential rulings (¶ 5 supra) were issued in a record in which the existence, provenance, and custody of the 911 audio, BWC footage, CAD logs, and supplemental report had been denied on the record by officers of the court. The ultimate Day-7/Day-8 leak-through of the 911 audio does not retroactively place the concealed record ``fairly before the court'' for the purpose of the seven rulings that had already issued; \emph{Modnick} fixes the measurement at the point of adjudication, not at the point of eventual admission.

34.6. The additional standard is \emph{Estate of Sanders} (1985) 40 Cal.3d 607, 616, and its Supreme Court source in \emph{Larrabee v. Tracy} (1943) 21 Cal.2d 645, 650--651: extrinsic fraud exists whenever the complained-of conduct ``effectively prevented \ldots{} a fair adversary hearing.'' The Court's decision-making matrix on MIL No.~1, the Jan.~27, 2025 categorical exclusion of the 911 transcripts, the Feb.~18, 2025 denial of continuance and \textsection{} 473 relief, the Feb.~28, 2025 burden-shift, and the RT 7 exclusion of the supplemental report were all produced inside a record distorted by the concealment. Under \emph{Sanders/Larrabee}, that distortion is itself the measure of extrinsic fraud.

\subsection{\textsection{}IV.5.5 --- Measure of Deprivation: Opportunity to Present, Not Evidentiary Admission}\label{iv.5.5-measure-of-deprivation-opportunity-to-present-not-evidentiary-admission}

34.7. The defense is expected to argue that the ultimate Day-7 admission of the 911 audio cures any prior concealment. That argument is foreclosed by four California Supreme Court holdings. \emph{Westphal v. Westphal} (1942) 20 Cal.2d 393, 397, fixes the metric at deprivation of ``opportunity to present.'' \emph{Kulchar v. Kulchar} (1969) 1 Cal.3d 467, 471, identifies the injury as denial of ``a fair adversary hearing.'' \emph{In re Marriage of Modnick} (1983) 33 Cal.3d 897, 907, requires that the concealed subject have been ``fairly before the court for adjudication.'' \emph{Estate of Sanders} (1985) 40 Cal.3d 607, 616 (quoting \emph{Larrabee v. Tracy} (1943) 21 Cal.2d 645, 650--651), measures the harm by conduct that ``effectively prevented \ldots{} a fair adversary hearing.'' Under each of these Supreme Court holdings, the test is whether Plaintiff was deprived of the opportunity to present his case on a non-corrupted record --- not whether one exhibit was admitted late under a corrupted posture. Plaintiff was deprived under each of the four Supreme Court tests. The Day-7 leak-through defense is legally irrelevant to the Supreme Court's deprivation metric.

\begin{center}\rule{0.5\linewidth}{0.5pt}\end{center}

\phantomsection\label{toc:section-vi-tiered-prayer-for-relief}
\section{SECTION VI: TIERED PRAYER FOR RELIEF}\label{section-vi-tiered-prayer-for-relief}

\subsection{Tier 1: Equitable Relief (Primary)}\label{tier-1-equitable-relief-primary}

\begin{enumerate}
\def\labelenumi{\arabic{enumi}.}
\setcounter{enumi}{34}
\tightlist
\item
  Plaintiff prays for the following equitable relief:
\end{enumerate}

\begin{enumerate}
\def\labelenumi{(\alph{enumi})}
\item
  \textbf{Vacatur:} An order vacating the judgment entered in CGC-21-594102 on April 25, 2025, on the ground that it was procured by extrinsic fraud.
\item
  \textbf{New Trial:} An order granting Plaintiff a new trial in CGC-21-594102, to be conducted without the taint of Defendants' fraudulent concealment.
\item
  \textbf{Injunctive Relief:} An order enjoining Defendants from enforcing any portion of the judgment in CGC-21-594102 pending resolution of this action.
\end{enumerate}

\subsection{Tier 2: Ancillary Restitutionary Relief (Fraud-Delta Costs Only)}\label{tier-2-ancillary-restitutionary-relief-fraud-delta-costs-only}

\begin{enumerate}
\def\labelenumi{\arabic{enumi}.}
\setcounter{enumi}{35}
\tightlist
\item
  In the alternative, or in addition to vacatur, Plaintiff prays for restitutionary relief limited to costs and expenses caused by Defendants' fraud:
\end{enumerate}

\begin{enumerate}
\def\labelenumi{(\alph{enumi})}
\item
  \textbf{Fraud-Delta Litigation Costs:} The incremental costs incurred by Plaintiff in litigating CGC-21-594102 that would not have been incurred but for Defendants' fraudulent concealment.
\item
  \textbf{Post-Judgment Discovery Costs:} Costs incurred in discovering the fraud after judgment, including costs of obtaining certified transcripts and reconstructing the contradiction chronology.
\end{enumerate}

\subsection{Tier 3: Statutory Damages}\label{tier-3-statutory-damages}

\begin{enumerate}
\def\labelenumi{\arabic{enumi}.}
\setcounter{enumi}{36}
\tightlist
\item
  Plaintiff prays for statutory damages and penalties as follows:
\end{enumerate}

\begin{enumerate}
\def\labelenumi{(\alph{enumi})}
\item
  \textbf{Business and Professions Code \textsection{} 6128(a)/(b):} Treble damages and other relief available under \textsection{} 6128 for attorney deceit and collusion. \emph{Rickley v. Goodfriend} (2013) 212 Cal.App.4th 1136.
\item
  \textbf{Code of Civil Procedure \textsection{} 128.5:} Sanctions against Defendants for bad faith actions or tactics that are frivolous or solely intended to cause unnecessary delay.
\item
  \textbf{Code of Civil Procedure \textsection{} 2023.030:} Discovery sanctions for misuse of the discovery process, including failure to respond to discovery and suppression of evidence.
\item
  \textbf{State Bar Referral under \textsection{} 6086.7:} Plaintiff requests the Court exercise its mandatory duty to report attorney conduct warranting discipline.
\end{enumerate}

\subsection{Tier 4: Prejudice Support (Emotional Distress as Equity Factor)}\label{tier-4-prejudice-support-emotional-distress-as-equity-factor}

\begin{enumerate}
\def\labelenumi{\arabic{enumi}.}
\setcounter{enumi}{37}
\tightlist
\item
  Plaintiff alleges that he has suffered severe emotional distress as a result of Defendants' fraudulent conduct. This allegation is made not as a separate cause of action for intentional infliction of emotional distress, but as a factor supporting the Court's exercise of equitable jurisdiction and as evidence of the prejudice caused by Defendants' fraud.
\end{enumerate}

\subsection{Alternative Terminating Sanctions Relief}\label{alternative-terminating-sanctions-relief}

\begin{enumerate}
\def\labelenumi{\arabic{enumi}.}
\setcounter{enumi}{38}
\tightlist
\item
  In the alternative to vacatur, Plaintiff requests terminating sanctions under \emph{Slesinger, Inc.~v. Walt Disney Co.} (2007) 155 Cal.App.4th 736, given that Defendants' corruption of the judicial process is so pervasive that no remedy short of terminating sanctions can restore the integrity of the case.
\end{enumerate}

\begin{center}\rule{0.5\linewidth}{0.5pt}\end{center}

\phantomsection\label{toc:section-vii-partyspecific-allegations}
\section{SECTION VII: PARTY-SPECIFIC ALLEGATIONS}\label{section-vii-party-specific-allegations}

\subsection{Defendant CSAA Insurance Exchange}\label{defendant-csaa-insurance-exchange}

\begin{enumerate}
\def\labelenumi{\arabic{enumi}.}
\setcounter{enumi}{39}
\item
  Defendant CSAA INSURANCE EXCHANGE is a California interinsurance exchange with its principal place of business at 3055 Oak Road, Mailstop W280, Walnut Creek, California 94597. CSAA was the liability insurer for defendant Subhi Abdelhalim in the underlying action CGC-21-594102 and, as institutional principal, retained, funded, directed, and ratified the conduct of defense counsel throughout that proceeding.
\item
  On February 25, 2021---just 21 days after the collision---CSAA issued a written denial letter for Claim No.~\textbf{1004-09-1054} under Policy No.~\textbf{CAAS100310857}, signed by \textbf{Sarah Rash}, Claims Representative, denying coverage before any meaningful investigation. This precipitous denial reflects CSAA's institutional posture throughout the litigation.
\end{enumerate}

\subsection{Defendant Carbone, Smith \& Koyama, LLP (CSK)}\label{defendant-carbone-smith-koyama-llp-csk}

\begin{enumerate}
\def\labelenumi{\arabic{enumi}.}
\setcounter{enumi}{41}
\tightlist
\item
  Defendant CARBONE, SMITH \& KOYAMA, LLP is a California limited liability partnership with offices at 555 12th Street, Suite 1250, Oakland, California 94607-4095. CSK served as defense counsel in CGC-21-594102 from approximately September 2021 through May 4, 2023, and had custody of the defense file---including discovery responses, deposition transcripts, and the 911/BWC/CAD materials---during that period.
\end{enumerate}

\subsection{Defendant Michael R. Chambers (State Bar No.~132394)}\label{defendant-michael-r.-chambers-state-bar-no.-132394}

\begin{enumerate}
\def\labelenumi{\arabic{enumi}.}
\setcounter{enumi}{42}
\tightlist
\item
  Defendant MICHAEL R. CHAMBERS is an individual attorney licensed to practice law in California and associated with CSK during the relevant period. Chambers engaged in the following specific conduct:
\end{enumerate}

\begin{enumerate}
\def\labelenumi{(\alph{enumi})}
\item
  Appeared at Plaintiff Rosario's August 18, 2022 deposition when the 911 audio was served on defense counsel;
\item
  Appeared at Defendant Abdelhalim's September 30, 2022 deposition when 911-related sworn testimony was given;
\item
  Executed the May 4, 2023 file transfer to successor counsel PSA without conveying the production history for the 911/BWC/CAD materials.
\end{enumerate}

\subsection{Defendant Phillips, Spallas \& Angstadt, LLP (PSA)}\label{defendant-phillips-spallas-angstadt-llp-psa}

\begin{enumerate}
\def\labelenumi{\arabic{enumi}.}
\setcounter{enumi}{43}
\tightlist
\item
  Defendant PHILLIPS, SPALLAS \& ANGSTADT, LLP is a California limited liability partnership with offices at 560 Mission Street, Suite 1010, San Francisco, California 94105. PSA substituted in as defense counsel on May 4, 2023, pursuant to the Substitution of Attorney (Transaction ID \#69953800), and represented defendant Abdelhalim through trial.
\end{enumerate}

\subsection{Defendant Priya D. Navaratnasingham (State Bar No.~252686)}\label{defendant-priya-d.-navaratnasingham-state-bar-no.-252686}

\begin{enumerate}
\def\labelenumi{\arabic{enumi}.}
\setcounter{enumi}{44}
\tightlist
\item
  Defendant PRIYA D. NAVARATNASINGHAM is an individual attorney licensed to practice law in California and associated with PSA. Navaratnasingham engaged in the following specific conduct:
\end{enumerate}

\begin{enumerate}
\def\labelenumi{(\alph{enumi})}
\item
  Signed and filed Motions in Limine Nos. 1--19 on January 17, 2025, including MIL No.~17 asserting ``Defendant has no idea where this 911 call log or audio recording came from'';
\item
  Made the ``news to me'' representation at the February 18, 2025 hearing (Hr'g Tr. p.~141, ll. 4--11): ``our office did not receive the 911 call until the day of the MSC\ldots{} this is news to me'';
\item
  Made the contradictory ``provided by Dolan'' admission on April 16, 2025 (RT 7, p.~213, ll. 13--15): ``Sure. They were provided by Dolan. He has had them for whatever, sure.''
\end{enumerate}

\subsection{Defendant Alberto Reyna (State Bar No.~347323)}\label{defendant-alberto-reyna-state-bar-no.-347323}

\begin{enumerate}
\def\labelenumi{\arabic{enumi}.}
\setcounter{enumi}{45}
\tightlist
\item
  Defendant ALBERTO REYNA is an individual attorney licensed to practice law in California and associated with PSA. Reyna engaged in the following specific conduct:
\end{enumerate}

\begin{enumerate}
\def\labelenumi{(\alph{enumi})}
\item
  Issued a Compex subpoena (Order No.~CA1147102-021) to SFPD on February 1, 2024, seeking the identical 911/BWC/CAD materials that were later portrayed as ``unknown'' or ``not received'';
\item
  Made representations at trial on April 16, 2025 that the supplemental police report was ``not anything that we received as part of the subpoena'' (RT 7, p.~156, ll. 5--16).
\end{enumerate}

\subsection{Defendants DOES 1-20}\label{defendants-does-1-20}

\begin{enumerate}
\def\labelenumi{\arabic{enumi}.}
\setcounter{enumi}{46}
\tightlist
\item
  Defendants DOES 1-20 are persons or entities whose true names and capacities are unknown to Plaintiff. Plaintiff will amend this complaint to state the true names and capacities when ascertained.
\end{enumerate}

\begin{center}\rule{0.5\linewidth}{0.5pt}\end{center}

\phantomsection\label{toc:section-viii-day7-admission-foundation}
\section{SECTION VIII: DAY-7 ADMISSION FOUNDATION}\label{section-viii-day-7-admission-foundation}

\subsection{Quoted Transcript --- April 16, 2025 (Trial Day 7)}\label{quoted-transcript-april-16-2025-trial-day-7}

\begin{enumerate}
\def\labelenumi{\arabic{enumi}.}
\setcounter{enumi}{47}
\tightlist
\item
  On April 16, 2025 (Trial Day 7 in CGC-21-594102), Defendant Navaratnasingham made the following admission on the certified record:
\end{enumerate}

\begin{quote}
``Sure. They were provided by Dolan. He has had them for whatever, sure.'' --- \emph{RT 7, p.~213, ll. 13--15.}
\end{quote}

\begin{enumerate}
\def\labelenumi{\arabic{enumi}.}
\setcounter{enumi}{48}
\tightlist
\item
  This admission---that the body-worn camera materials came from Dolan Law Firm (the same source as the 911 audio in Jessup's testimony)---is irreconcilable with:
\end{enumerate}

\begin{enumerate}
\def\labelenumi{(\alph{enumi})}
\item
  Navaratnasingham's February 18, 2025 ``news to me'' statement;
\item
  MIL No.~17's assertion that ``Defendant has no idea where this 911 call log or audio recording came from''; and
\item
  The January 27, 2025 imputation of fabrication: ``we really don't know who redacted these things and what --- if that representation is true'' (Hr'g Tr. Jan.~27, 2025, p.~31, ll. 1--5).
\end{enumerate}

\begin{enumerate}
\def\labelenumi{\arabic{enumi}.}
\setcounter{enumi}{49}
\tightlist
\item
  Despite this admission, defense counsel did not:
\end{enumerate}

\begin{enumerate}
\def\labelenumi{(\alph{enumi})}
\item
  Correct the record regarding the source of the materials;
\item
  Acknowledge the inconsistency with prior tribunal-facing representations;
\item
  Withdraw or modify MIL Nos. 6 or 17; or
\item
  Inform the Court that the prior ``never received'' representations were false.
\end{enumerate}

\begin{center}\rule{0.5\linewidth}{0.5pt}\end{center}

\phantomsection\label{toc:section-ix-lawsbroken-catalog}
\section{SECTION IX: LAWS-BROKEN CATALOG}\label{section-ix-laws-broken-catalog}

\begin{enumerate}
\def\labelenumi{\arabic{enumi}.}
\setcounter{enumi}{50}
\tightlist
\item
  The following laws were violated by Defendants' conduct:
\end{enumerate}

{\def\LTcaptype{table} % do not increment counter
\begin{longtable}[]{@{}
  >{\raggedright\arraybackslash}p{(\linewidth - 4\tabcolsep) * \real{0.3333}}
  >{\raggedright\arraybackslash}p{(\linewidth - 4\tabcolsep) * \real{0.3333}}
  >{\raggedright\arraybackslash}p{(\linewidth - 4\tabcolsep) * \real{0.3333}}@{}}
\toprule\noalign{}
\begin{minipage}[b]{\linewidth}\raggedright
Rule / Code
\end{minipage} & \begin{minipage}[b]{\linewidth}\raggedright
Provision
\end{minipage} & \begin{minipage}[b]{\linewidth}\raggedright
Defense Conduct
\end{minipage} \\
\midrule\noalign{}
\endhead
\bottomrule\noalign{}
\endlastfoot
Bus. \& Prof.~Code \textsection{} 6068(d) & Duty not to mislead judge & Jan 27, Feb 18, Feb 28 denials re receipt \\
Bus. \& Prof.~Code \textsection{} 6128(a) & Attorney deceit / collusion & Intentional deception re: file transfer and receipt \\
Cal. Rule of Prof.~Conduct 3.3(a) & Candor to tribunal & ``We never received these CDs, certainly'' (Jan 27) \\
Cal. Rule of Prof.~Conduct 3.4(a) & Fairness to opposing counsel & Withholding receipt history to procure exclusion \\
Cal. Rule of Prof.~Conduct 8.4(c) & Dishonesty / Misrepresentation & Sustained pattern across four hearings \\
CCP \textsection{} 128.5 & Bad-faith actions & Procurement of exclusion through false statements \\
Civ. Code \textsection{} 1572(3) & Actual fraud (suppression) & Concealment of receipt for 893 days \\
Civ. Code \textsection{} 1573 & Constructive fraud & Breach of duty of candor; misleading acts \\
Evid. Code \textsection{} 413 & Willful suppression & Adverse inference applicable to 911-related issues \\
\end{longtable}
}

\begin{center}\rule{0.5\linewidth}{0.5pt}\end{center}

\phantomsection\label{toc:causes-of-action}
\section{CAUSES OF ACTION}\label{causes-of-action}

\subsection{FIRST CAUSE OF ACTION: Extrinsic Fraud --- Independent Action in Equity (Against All Defendants)}\label{first-cause-of-action-extrinsic-fraud-independent-action-in-equity-against-all-defendants}

\begin{enumerate}
\def\labelenumi{\arabic{enumi}.}
\setcounter{enumi}{51}
\item
  Plaintiff incorporates by reference all preceding paragraphs.
\item
  Defendants, by the conduct alleged herein, procured the judgment in CGC-21-594102 through extrinsic fraud operating outside the trial itself.
\item
  The fraud deprived Plaintiff of a full and fair opportunity to present his case by concealing the production history of critical evidence and inducing the Court to issue structural exclusionary rulings.
\item
  Plaintiff is entitled to vacatur of the judgment, a new trial, and such other relief as the Court deems equitable.
\end{enumerate}

\subsection{SECOND CAUSE OF ACTION: Constructive Fraud --- Civil Code \textsection{} 1573 (Against All Defendants)}\label{second-cause-of-action-constructive-fraud-civil-code-1573-against-all-defendants}

\begin{enumerate}
\def\labelenumi{\arabic{enumi}.}
\setcounter{enumi}{55}
\item
  Plaintiff incorporates by reference all preceding paragraphs.
\item
  Defendants owed Plaintiff a duty not to gain an unfair advantage through concealment or deception, particularly in the context of litigation where the Court's processes depend on candor and fair dealing.
\item
  The February 28, 2025 colloquy establishes that the tribunal extended confidential institutional reliance to counsel's representations about possession and provenance.
\item
  Defendants breached this duty by concealing the production history, making materially false representations about receipt, and failing to correct the record after the April 16, 2025 admission.
\item
  As a direct and proximate result of Defendants' constructive fraud, Plaintiff has suffered damages as alleged herein.
\end{enumerate}

\subsection{THIRD CAUSE OF ACTION: Violation of Business and Professions Code \textsection{} 6128(a) (Against Attorney Defendants)}\label{third-cause-of-action-violation-of-business-and-professions-code-6128a-against-attorney-defendants}

\begin{enumerate}
\def\labelenumi{\arabic{enumi}.}
\setcounter{enumi}{60}
\item
  Plaintiff incorporates by reference all preceding paragraphs.
\item
  Defendants Chambers, Navaratnasingham, Reyna, CSK, and PSA, by the conduct alleged herein, engaged in deceit or collusion, or consented to deceit or collusion, with intent to deceive the court or Plaintiff, in violation of Business and Professions Code \textsection{} 6128(a).
\item
  As a direct and proximate result of this violation, Plaintiff has suffered damages as alleged herein, and is entitled to treble damages under \textsection{} 6128(b). \emph{Rickley v. Goodfriend} (2013) 212 Cal.App.4th 1136.
\end{enumerate}

\subsection{FOURTH CAUSE OF ACTION: Negligence Per Se (Against All Defendants)}\label{fourth-cause-of-action-negligence-per-se-against-all-defendants}

\begin{enumerate}
\def\labelenumi{\arabic{enumi}.}
\setcounter{enumi}{63}
\item
  Plaintiff incorporates by reference all preceding paragraphs.
\item
  Defendants violated Business and Professions Code \textsection{} 6128(a), a statute designed to protect litigants from attorney deceit and collusion.
\item
  This violation constitutes negligence per se, and Plaintiff is entitled to damages caused by this negligence.
\end{enumerate}

\subsection{FIFTH CAUSE OF ACTION: Fraud --- Intentional Misrepresentation (Against Attorney Defendants)}\label{fifth-cause-of-action-fraud-intentional-misrepresentation-against-attorney-defendants}

\begin{enumerate}
\def\labelenumi{\arabic{enumi}.}
\setcounter{enumi}{66}
\item
  Plaintiff incorporates by reference all preceding paragraphs.
\item
  Defendants made materially false representations to the tribunal regarding their receipt, custody, and knowledge of the 911/BWC/CAD materials.
\item
  Defendants knew these representations were false at the time they were made, as evidenced by the April 16, 2025 admission.
\item
  The Court relied on these representations in issuing the seven consequential rulings identified in ¶ 5.
\item
  As a direct and proximate result of Defendants' intentional misrepresentations, Plaintiff suffered harm including entry of an adverse judgment, deprivation of a fair trial, and the costs of prosecuting this action.
\end{enumerate}

\subsection{SIXTH CAUSE OF ACTION: Aiding and Abetting Fraud (Against CSAA)}\label{sixth-cause-of-action-aiding-and-abetting-fraud-against-csaa}

\begin{enumerate}
\def\labelenumi{\arabic{enumi}.}
\setcounter{enumi}{71}
\item
  Plaintiff incorporates by reference all preceding paragraphs.
\item
  CSAA, as the institutional principal funding, directing, and ratifying defense counsel's litigation conduct, aided and abetted the fraud perpetrated by its panel counsel.
\item
  Under \emph{Doctors' Company v. Superior Court} (1989) 49 Cal.3d 39, an insurer that acts outside its normal insurer role may be held liable for aiding and abetting its counsel's fraudulent conduct.
\item
  As a direct and proximate result of CSAA's conduct, Plaintiff suffered damages as alleged herein.
\end{enumerate}

\subsection{SEVENTH CAUSE OF ACTION: Constructive Fraud --- Civil Code \textsection{} 1573 (Against Attorney Defendants)}\label{seventh-cause-of-action-constructive-fraud-civil-code-1573-against-attorney-defendants}

\begin{enumerate}
\def\labelenumi{\arabic{enumi}.}
\setcounter{enumi}{75}
\item
  Plaintiff incorporates by reference all preceding paragraphs.
\item
  The individual attorney defendants owed Plaintiff, as a party to the litigation, a duty of candor to the tribunal that inured to Plaintiff's benefit as the opposing party.
\item
  Defendants breached this duty by concealing the production history and making representations designed to mislead the tribunal about their possession of critical evidence.
\item
  Under \emph{Assilzadeh v. California Federal Bank} (2000) 82 Cal.App.4th 399, 415, constructive fraud under Civil Code \textsection{} 1573 does not require proof of actual fraudulent intent---breach of a legal or equitable duty arising from a confidential or fiduciary relationship, by a misleading act or nondisclosure, is sufficient.
\item
  As a direct and proximate result of Defendants' constructive fraud, Plaintiff suffered damages as alleged herein.
\end{enumerate}

\subsection{EIGHTH CAUSE OF ACTION: Extrinsic Fraud by Breach of Officer-of-the-Court Fiduciary Duty (Against Attorney Defendants)}\label{eighth-cause-of-action-extrinsic-fraud-by-breach-of-officer-of-the-court-fiduciary-duty-against-attorney-defendants}

\begin{enumerate}
\def\labelenumi{\arabic{enumi}.}
\setcounter{enumi}{80}
\item
  Plaintiff incorporates by reference all preceding paragraphs.
\item
  Defendants Chambers, Navaratnasingham, Reyna, CSK, and PSA are officers of the court. The California Supreme Court in \emph{Estate of Sanders} (1985) 40 Cal.3d 607, 615, holds that officer-of-the-court status imposes a fiduciary duty running to the integrity of the adjudicative process and, derivatively, to the opposing party's primary right to a fair adversary hearing.
\item
  The February 28, 2025 colloquy (Hr'g Tr. pp.~291:26--292:7) is the on-record predicate that the trial court extended institutional reliance to this fiduciary status. The Court stated: ``I'm not going to question a lawyer's remarks regarding what they have and don't have. They are an officer of the Court, and they are supposed to conduct themselves appropriately before the Court.''
\item
  The Attorney Defendants breached that fiduciary duty by:
\end{enumerate}

\begin{enumerate}
\def\labelenumi{(\alph{enumi})}
\item
  Maintaining categorical denials of receipt of materials they had possessed since August-September 2022;
\item
  Executing the May 4, 2023 file transfer without conveying the production history;
\item
  Procuring structural rulings --- MIL No.~1, the Jan.~27, 2025 categorical exclusion of the 911 transcripts, the Feb.~18, 2025 denial of continuance and \textsection{} 473 relief, the Feb.~28, 2025 burden-shift, and the RT 7 exclusion of the supplemental report --- by representations that the April 16, 2025 admission later exposed as false; and
\item
  Failing to correct the record after the April 16, 2025 admission that the materials ``were provided by Dolan.''
\end{enumerate}

\begin{enumerate}
\def\labelenumi{\arabic{enumi}.}
\setcounter{enumi}{84}
\item
  Under \emph{Estate of Sanders}, 40 Cal.3d at 615--616, this fiduciary breach by officers of the court that ``effectively prevented'' a fair adversary hearing is, as a matter of California Supreme Court doctrine, extrinsic fraud. Under \emph{In re Marriage of Modnick} (1983) 33 Cal.3d 897, 907, that conduct produced a record in which the concealed evidentiary subject was ``never fairly before the court for adjudication.''
\item
  This cause of action is additive to, and does not displace, the independent action in equity pleaded in the First Cause of Action. It is pleaded separately to state the officer-of-the-court fiduciary-duty theory recognized by the Supreme Court in \emph{Sanders}, in discrete form, against the Attorney Defendants.
\item
  As a direct and proximate result of the Attorney Defendants' breach of officer-of-the-court fiduciary duty, Plaintiff was deprived of a fair adversary hearing within the meaning of \emph{Westphal} (1942) 20 Cal.2d 393, 397; \emph{Kulchar} (1969) 1 Cal.3d 467, 471; \emph{Modnick} (1983) 33 Cal.3d 897, 905, 907; and \emph{Sanders} (1985) 40 Cal.3d 607, 614--616. Plaintiff is entitled to the equitable relief sought in this action, together with such other and further relief as the Court deems just and proper.
\end{enumerate}

\begin{center}\rule{0.5\linewidth}{0.5pt}\end{center}

\phantomsection\label{toc:prayer-for-relief}
\section{PRAYER FOR RELIEF}\label{prayer-for-relief}

WHEREFORE, Plaintiff prays for judgment against Defendants as follows:

\textbf{Equitable Relief:} 1. Vacatur of the judgment entered in CGC-21-594102 on April 25, 2025; 2. A new trial in CGC-21-594102; 3. Preliminary and permanent injunctive relief; 4. Terminating sanctions under \emph{Slesinger} in the alternative;

\textbf{Statutory Damages:} 5. Treble damages under Business and Professions Code \textsection{} 6128(b); 6. Sanctions under Code of Civil Procedure \textsection{}\textsection{} 128.5 and 2023.030; 7. State Bar referral under Business and Professions Code \textsection{} 6086.7;

\textbf{Compensatory Damages:} 8. Restitutionary damages limited to fraud-delta costs; 9. Post-judgment discovery and reconstruction costs;

\textbf{Costs and General Relief:} 10. Costs of suit incurred herein; 11. Such other and further relief as the Court deems just and proper.

\begin{center}\rule{0.5\linewidth}{0.5pt}\end{center}

\textbf{Dated:} \_\_\_\_\_\_\_\_\_\_\_\_\_\_\_

\textbf{FRANCISCUS DYLAN ROSARIO}\\
Plaintiff, In Pro Per\\
7624 Melody Drive\\
Rohnert Park, CA 94928\\
Tel: 650.488.7510\\
Email: soltrinox@gmail.com

\begin{center}\rule{0.5\linewidth}{0.5pt}\end{center}

\phantomsection\label{toc:verification}
\section{VERIFICATION}\label{verification}

I, FRANCISCUS DYLAN ROSARIO, am the Plaintiff in this action. I have read the foregoing Third Amended Complaint and know its contents. The facts stated therein are true of my own knowledge, except as to matters stated on information and belief, and as to those matters, I believe them to be true.

I declare under penalty of perjury under the laws of the State of California that the foregoing is true and correct.

Executed on \_\_\_\_\_\_\_\_\_\_\_\_\_\_\_ at \_\_\_\_\_\_\_\_\_\_\_\_\_\_\_, California.

\begin{center}\rule{0.5\linewidth}{0.5pt}\end{center}

FRANCISCUS DYLAN ROSARIO

\nolinenumbers
\vspace{6mm}
\noindent /s/ Franciscus Dylan Rosario\\
\vspace{4mm}
\noindent\includegraphics[height=0.5in]{../../../Support/signature.png}

\end{document}
